\PassOptionsToPackage{unicode=true}{hyperref} % options for packages loaded elsewhere
\PassOptionsToPackage{hyphens}{url}
%
\documentclass[ignorenonframetext,]{beamer}
\usepackage{pgfpages}
\setbeamertemplate{caption}[numbered]
\setbeamertemplate{caption label separator}{: }
\setbeamercolor{caption name}{fg=normal text.fg}
\beamertemplatenavigationsymbolsempty
\usepackage{lmodern}
\usepackage{amssymb,amsmath}
\usepackage{ifxetex,ifluatex}
\usepackage{fixltx2e} % provides \textsubscript
\ifnum 0\ifxetex 1\fi\ifluatex 1\fi=0 % if pdftex
  \usepackage[T1]{fontenc}
  \usepackage[utf8]{inputenc}
  \usepackage{textcomp} % provides euro and other symbols
\else % if luatex or xelatex
  \usepackage{unicode-math}
  \defaultfontfeatures{Ligatures=TeX,Scale=MatchLowercase}
\fi
\usetheme[]{CambridgeUS}
\usecolortheme{beaver}
\usefonttheme{structurebold}
% use upquote if available, for straight quotes in verbatim environments
\IfFileExists{upquote.sty}{\usepackage{upquote}}{}
% use microtype if available
\IfFileExists{microtype.sty}{%
\usepackage[]{microtype}
\UseMicrotypeSet[protrusion]{basicmath} % disable protrusion for tt fonts
}{}
\IfFileExists{parskip.sty}{%
\usepackage{parskip}
}{% else
\setlength{\parindent}{0pt}
\setlength{\parskip}{6pt plus 2pt minus 1pt}
}
\usepackage{hyperref}
\hypersetup{
            pdftitle={A8 Rasterdaten},
            pdfauthor={Jan-Philipp Kolb},
            pdfborder={0 0 0},
            breaklinks=true}
\urlstyle{same}  % don't use monospace font for urls
\newif\ifbibliography
\usepackage{color}
\usepackage{fancyvrb}
\newcommand{\VerbBar}{|}
\newcommand{\VERB}{\Verb[commandchars=\\\{\}]}
\DefineVerbatimEnvironment{Highlighting}{Verbatim}{commandchars=\\\{\}}
% Add ',fontsize=\small' for more characters per line
\usepackage{framed}
\definecolor{shadecolor}{RGB}{42,33,28}
\newenvironment{Shaded}{\begin{snugshade}}{\end{snugshade}}
\newcommand{\AlertTok}[1]{\textcolor[rgb]{1.00,1.00,0.00}{#1}}
\newcommand{\AnnotationTok}[1]{\textcolor[rgb]{0.00,0.40,1.00}{\textbf{\textit{#1}}}}
\newcommand{\AttributeTok}[1]{\textcolor[rgb]{0.74,0.68,0.62}{#1}}
\newcommand{\BaseNTok}[1]{\textcolor[rgb]{0.27,0.67,0.26}{#1}}
\newcommand{\BuiltInTok}[1]{\textcolor[rgb]{0.74,0.68,0.62}{#1}}
\newcommand{\CharTok}[1]{\textcolor[rgb]{0.02,0.61,0.04}{#1}}
\newcommand{\CommentTok}[1]{\textcolor[rgb]{0.00,0.40,1.00}{\textbf{\textit{#1}}}}
\newcommand{\CommentVarTok}[1]{\textcolor[rgb]{0.74,0.68,0.62}{#1}}
\newcommand{\ConstantTok}[1]{\textcolor[rgb]{0.74,0.68,0.62}{#1}}
\newcommand{\ControlFlowTok}[1]{\textcolor[rgb]{0.26,0.66,0.93}{\textbf{#1}}}
\newcommand{\DataTypeTok}[1]{\textcolor[rgb]{0.74,0.68,0.62}{\underline{#1}}}
\newcommand{\DecValTok}[1]{\textcolor[rgb]{0.27,0.67,0.26}{#1}}
\newcommand{\DocumentationTok}[1]{\textcolor[rgb]{0.00,0.40,1.00}{\textit{#1}}}
\newcommand{\ErrorTok}[1]{\textcolor[rgb]{1.00,1.00,0.00}{\textbf{#1}}}
\newcommand{\ExtensionTok}[1]{\textcolor[rgb]{0.74,0.68,0.62}{#1}}
\newcommand{\FloatTok}[1]{\textcolor[rgb]{0.27,0.67,0.26}{#1}}
\newcommand{\FunctionTok}[1]{\textcolor[rgb]{1.00,0.58,0.35}{\textbf{#1}}}
\newcommand{\ImportTok}[1]{\textcolor[rgb]{0.74,0.68,0.62}{#1}}
\newcommand{\InformationTok}[1]{\textcolor[rgb]{0.00,0.40,1.00}{\textbf{\textit{#1}}}}
\newcommand{\KeywordTok}[1]{\textcolor[rgb]{0.26,0.66,0.93}{\textbf{#1}}}
\newcommand{\NormalTok}[1]{\textcolor[rgb]{0.74,0.68,0.62}{#1}}
\newcommand{\OperatorTok}[1]{\textcolor[rgb]{0.74,0.68,0.62}{#1}}
\newcommand{\OtherTok}[1]{\textcolor[rgb]{0.74,0.68,0.62}{#1}}
\newcommand{\PreprocessorTok}[1]{\textcolor[rgb]{0.74,0.68,0.62}{\textbf{#1}}}
\newcommand{\RegionMarkerTok}[1]{\textcolor[rgb]{0.74,0.68,0.62}{#1}}
\newcommand{\SpecialCharTok}[1]{\textcolor[rgb]{0.02,0.61,0.04}{#1}}
\newcommand{\SpecialStringTok}[1]{\textcolor[rgb]{0.02,0.61,0.04}{#1}}
\newcommand{\StringTok}[1]{\textcolor[rgb]{0.02,0.61,0.04}{#1}}
\newcommand{\VariableTok}[1]{\textcolor[rgb]{0.74,0.68,0.62}{#1}}
\newcommand{\VerbatimStringTok}[1]{\textcolor[rgb]{0.02,0.61,0.04}{#1}}
\newcommand{\WarningTok}[1]{\textcolor[rgb]{1.00,1.00,0.00}{\textbf{#1}}}
\usepackage{graphicx,grffile}
\makeatletter
\def\maxwidth{\ifdim\Gin@nat@width>\linewidth\linewidth\else\Gin@nat@width\fi}
\def\maxheight{\ifdim\Gin@nat@height>\textheight\textheight\else\Gin@nat@height\fi}
\makeatother
% Scale images if necessary, so that they will not overflow the page
% margins by default, and it is still possible to overwrite the defaults
% using explicit options in \includegraphics[width, height, ...]{}
\setkeys{Gin}{width=\maxwidth,height=\maxheight,keepaspectratio}
% Prevent slide breaks in the middle of a paragraph:
\widowpenalties 1 10000
\raggedbottom
\setbeamertemplate{part page}{
\centering
\begin{beamercolorbox}[sep=16pt,center]{part title}
  \usebeamerfont{part title}\insertpart\par
\end{beamercolorbox}
}
\setbeamertemplate{section page}{
\centering
\begin{beamercolorbox}[sep=12pt,center]{part title}
  \usebeamerfont{section title}\insertsection\par
\end{beamercolorbox}
}
\setbeamertemplate{subsection page}{
\centering
\begin{beamercolorbox}[sep=8pt,center]{part title}
  \usebeamerfont{subsection title}\insertsubsection\par
\end{beamercolorbox}
}
\AtBeginPart{
  \frame{\partpage}
}
\AtBeginSection{
  \ifbibliography
  \else
    \frame{\sectionpage}
  \fi
}
\AtBeginSubsection{
  \frame{\subsectionpage}
}
\setlength{\emergencystretch}{3em}  % prevent overfull lines
\providecommand{\tightlist}{%
  \setlength{\itemsep}{0pt}\setlength{\parskip}{0pt}}
\setcounter{secnumdepth}{0}

% set default figure placement to htbp
\makeatletter
\def\fps@figure{htbp}
\makeatother


\title{A8 Rasterdaten}
\author{Jan-Philipp Kolb}
\date{22 Oktober 2018}

\begin{document}
\frame{\titlepage}

\begin{frame}[fragile]{\href{http://rspatial.org/spatial/rst/8-rastermanip.html}{Manipulation
von Rasterdaten}}
\protect\hypertarget{manipulation-von-rasterdaten}{}

\begin{Shaded}
\begin{Highlighting}[]
\KeywordTok{library}\NormalTok{(raster)}
\end{Highlighting}
\end{Shaded}

\begin{verbatim}
## Loading required package: sp
\end{verbatim}

\begin{Shaded}
\begin{Highlighting}[]
\NormalTok{x <-}\StringTok{ }\KeywordTok{raster}\NormalTok{()}
\NormalTok{x}
\end{Highlighting}
\end{Shaded}

\begin{verbatim}
## class       : RasterLayer 
## dimensions  : 180, 360, 64800  (nrow, ncol, ncell)
## resolution  : 1, 1  (x, y)
## extent      : -180, 180, -90, 90  (xmin, xmax, ymin, ymax)
## coord. ref. : +proj=longlat +datum=WGS84 +ellps=WGS84 +towgs84=0,0,0
\end{verbatim}

\end{frame}

\begin{frame}[fragile]{Andere Parameter wählen}
\protect\hypertarget{andere-parameter-wahlen}{}

\begin{Shaded}
\begin{Highlighting}[]
\NormalTok{x1 <-}\StringTok{ }\KeywordTok{raster}\NormalTok{(}\DataTypeTok{ncol=}\DecValTok{36}\NormalTok{, }\DataTypeTok{nrow=}\DecValTok{18}\NormalTok{, }\DataTypeTok{xmn=}\OperatorTok{-}\DecValTok{1000}\NormalTok{, }\DataTypeTok{xmx=}\DecValTok{1000}\NormalTok{, }\DataTypeTok{ymn=}\OperatorTok{-}\DecValTok{100}\NormalTok{, }\DataTypeTok{ymx=}\DecValTok{900}\NormalTok{)}
\NormalTok{x1}
\end{Highlighting}
\end{Shaded}

\begin{verbatim}
## class       : RasterLayer 
## dimensions  : 18, 36, 648  (nrow, ncol, ncell)
## resolution  : 55.55556, 55.55556  (x, y)
## extent      : -1000, 1000, -100, 900  (xmin, xmax, ymin, ymax)
## coord. ref. : NA
\end{verbatim}

\end{frame}

\begin{frame}[fragile]{Den Zellen Werte zuordnen}
\protect\hypertarget{den-zellen-werte-zuordnen}{}

\begin{Shaded}
\begin{Highlighting}[]
\NormalTok{r <-}\StringTok{ }\KeywordTok{raster}\NormalTok{(}\DataTypeTok{ncol=}\DecValTok{10}\NormalTok{, }\DataTypeTok{nrow=}\DecValTok{10}\NormalTok{)}
\KeywordTok{ncell}\NormalTok{(r)}
\end{Highlighting}
\end{Shaded}

\begin{verbatim}
## [1] 100
\end{verbatim}

\begin{Shaded}
\begin{Highlighting}[]
\KeywordTok{hasValues}\NormalTok{(r)}
\end{Highlighting}
\end{Shaded}

\begin{verbatim}
## [1] FALSE
\end{verbatim}

\begin{Shaded}
\begin{Highlighting}[]
\KeywordTok{values}\NormalTok{(r) <-}\StringTok{ }\KeywordTok{runif}\NormalTok{(}\KeywordTok{ncell}\NormalTok{(r))}
\KeywordTok{hasValues}\NormalTok{(r)}
\end{Highlighting}
\end{Shaded}

\begin{verbatim}
## [1] TRUE
\end{verbatim}

\end{frame}

\begin{frame}[fragile]{Das Ergebnis visualisieren}
\protect\hypertarget{das-ergebnis-visualisieren}{}

\begin{Shaded}
\begin{Highlighting}[]
\KeywordTok{plot}\NormalTok{(r, }\DataTypeTok{main=}\StringTok{'Raster with 100 cells'}\NormalTok{)}
\end{Highlighting}
\end{Shaded}

\includegraphics{A8_Rasterdaten_files/figure-beamer/unnamed-chunk-5-1.pdf}

--\textgreater{}

\end{frame}

\begin{frame}[fragile]{Import von Beispieldaten}
\protect\hypertarget{import-von-beispieldaten}{}

\begin{Shaded}
\begin{Highlighting}[]
\NormalTok{link<-}\StringTok{"https://raw.githubusercontent.com/GeoScripting-WUR/IntroToRaster/gh-pages/data/gewata.zip"}
\end{Highlighting}
\end{Shaded}

\begin{Shaded}
\begin{Highlighting}[]
\KeywordTok{download.file}\NormalTok{(}\DataTypeTok{url =}\NormalTok{ link, }\DataTypeTok{destfile =} \StringTok{'gewata.zip'}\NormalTok{, }\DataTypeTok{method =} \StringTok{'auto'}\NormalTok{)}
\KeywordTok{unzip}\NormalTok{(}\StringTok{'gewata.zip'}\NormalTok{)}
\end{Highlighting}
\end{Shaded}

\begin{Shaded}
\begin{Highlighting}[]
\KeywordTok{list.files}\NormalTok{(}\DataTypeTok{pattern =} \KeywordTok{glob2rx}\NormalTok{(}\StringTok{'*.tif'}\NormalTok{))}
\end{Highlighting}
\end{Shaded}

\begin{verbatim}
## character(0)
\end{verbatim}

\end{frame}

\begin{frame}[fragile]{\href{https://geoscripting-wur.github.io/IntroToRaster/}{Import
von Rasterdaten - \texttt{.tif} file}}
\protect\hypertarget{import-von-rasterdaten---.tif-file}{}

\begin{Shaded}
\begin{Highlighting}[]
\NormalTok{gewata <-}\StringTok{ }\NormalTok{raster}\OperatorTok{::}\KeywordTok{brick}\NormalTok{(}\StringTok{'../data/}
\StringTok{LE71700552001036SGS00_SR_Gewata_INT1U.tif'}\NormalTok{)}
\end{Highlighting}
\end{Shaded}

\end{frame}

\begin{frame}[fragile]{Erste plots erstellen}
\protect\hypertarget{erste-plots-erstellen}{}

\begin{Shaded}
\begin{Highlighting}[]
\KeywordTok{plot}\NormalTok{(gewata)}
\end{Highlighting}
\end{Shaded}

\includegraphics{A8_Rasterdaten_files/figure-beamer/unnamed-chunk-15-1.pdf}

\end{frame}

\begin{frame}[fragile]{\href{https://pakillo.github.io/R-GIS-tutorial/\#raster}{Worldclim
Daten importieren}}
\protect\hypertarget{worldclim-daten-importieren}{}

\includegraphics{figure/WORLDCLIM.PNG}

\begin{Shaded}
\begin{Highlighting}[]
\CommentTok{# this will download }
\CommentTok{# global data on minimum temperature at 10' resolution}
\NormalTok{tmin <-}\StringTok{ }\NormalTok{raster}\OperatorTok{::}\KeywordTok{getData}\NormalTok{(}\StringTok{"worldclim"}\NormalTok{, }\DataTypeTok{var =} \StringTok{"tmin"}\NormalTok{, }\DataTypeTok{res =} \DecValTok{10}\NormalTok{)}
\end{Highlighting}
\end{Shaded}

\begin{Shaded}
\begin{Highlighting}[]
\NormalTok{tmin1 <-}\StringTok{ }\NormalTok{raster}\OperatorTok{::}\KeywordTok{raster}\NormalTok{( }\StringTok{"../data/wc10/tmin1.bil"}\NormalTok{)}
\NormalTok{tmin1 <-}\StringTok{ }\NormalTok{tmin1}\OperatorTok{/}\DecValTok{10}
\end{Highlighting}
\end{Shaded}

\end{frame}

\begin{frame}[fragile]{Minimum Temperatur plotten}
\protect\hypertarget{minimum-temperatur-plotten}{}

\begin{Shaded}
\begin{Highlighting}[]
\KeywordTok{plot}\NormalTok{(tmin1)}
\end{Highlighting}
\end{Shaded}

\includegraphics{A8_Rasterdaten_files/figure-beamer/unnamed-chunk-19-1.pdf}

\end{frame}

\begin{frame}{Links}
\protect\hypertarget{links}{}

\begin{itemize}
\tightlist
\item
  Neon -
  \href{https://www.neonscience.org/dc-raster-data-r}{\textbf{Intro to
  Raster Data in R}}
\end{itemize}

\end{frame}

\end{document}
